\chapter*{Package guide}
\addcontentsline{toc}{chapter}{Package guide}

Load a package to solve a named requirement. Check its current manual, class
compatibility, engine support, licence, and build implications. A package listed
here is a useful category example, not a universal recommendation.

\begin{longtable}{@{}p{.22\textwidth}p{.32\textwidth}p{.38\textwidth}@{}}
\toprule
Package or tool & Main purpose & Use and cautions \\
\midrule
\endhead
\pkg{amsmath} &
Structured displayed mathematics. &
Provides \code{align}, \code{gather}, \code{cases}, \cmd{text}, and robust
equation tools. Prefer it to legacy display environments. \\
\addlinespace
\pkg{amssymb} and \pkg{amsfonts} &
Additional mathematical symbols and fonts. &
Load when their symbol sets are used; verify that a publisher class does not
already provide them. \\
\addlinespace
\pkg{mathtools} &
Extensions and refinements to \pkg{amsmath}. &
Useful for paired delimiters, specialised alignment, and mathematical tools.
Read its manual before adopting convenience commands. \\
\addlinespace
\pkg{amsthm} &
Theorem-like environments and proof styling. &
Define theorem structures once and follow the mathematical venue's numbering
convention. Some classes provide their own theorem interface. \\
\addlinespace
\pkg{siunitx} &
Numbers, quantities, units, and numeric table columns. &
Use \cmd{num}, \cmd{qty}, \cmd{unit}, and \code{S}. Version 3 syntax differs
from older examples; consult the installed manual. \\
\addlinespace
\pkg{graphicx} &
Insert and transform external graphics. &
Use project-relative paths and meaningful placed size. Keep editable or
generated figure sources and permission records. \\
\addlinespace
\pkg{caption} &
Caption formatting and options. &
Generic classes benefit from its interface, but publisher classes may control
captions and conflict with local restyling. \\
\addlinespace
\pkg{subcaption} &
Subfigures and subtables. &
Provides panel environments and labels. It builds on \pkg{caption}; verify class
compatibility and avoid the obsolete \pkg{subfigure} package. \\
\addlinespace
\pkg{booktabs} &
Restrained professional table rules. &
Use top, middle, and bottom rules; avoid vertical rules and double-rule grids
unless a domain-specific requirement says otherwise. \\
\addlinespace
\pkg{array} &
Extended column definitions. &
Useful for custom alignment and paragraph-column modifiers. Keep definitions
readable and local to a documented need. \\
\addlinespace
\pkg{tabularx} &
Tables with columns that fill a target width. &
The \code{X} column wraps prose. Combine cautiously with numeric columns and
avoid over-constrained width specifications. \\
\addlinespace
\pkg{longtable} &
Tables that break across pages. &
Runs directly in the page stream rather than as an ordinary float. Define
repeating headers and continuation rules. \\
\addlinespace
\pkg{xcolor} &
Named colours and colour models. &
Use colour sparingly, maintain contrast and grayscale meaning, and obey venue
rules. Colour cannot be the only carrier of information. \\
\addlinespace
\pkg{microtype} &
Subtle typographic protrusion and font expansion. &
Often improves justified text under pdf\LaTeX{}. Support depends on engine and
font; inspect rather than expecting a visible effect. \\
\addlinespace
\pkg{geometry} &
Explicit page dimensions and margins. &
Appropriate for generic documents with known requirements. Do not override
publisher or institutional page design without permission. \\
\addlinespace
\pkg{fancyhdr} &
Running headers and footers. &
Clear inherited fields, set a sufficient head height, and define a separate
plain style for opening pages when needed. \\
\addlinespace
\pkg{titlesec} &
Heading format and spacing in generic classes. &
Useful for a controlled house design; can conflict with specialist classes and
should not be used to imitate publisher formatting. \\
\addlinespace
\pkg{setspace} &
Single, one-and-a-half, and double spacing. &
Use when a verified review or thesis requirement asks for it. It handles more
contexts than guessing a baseline-stretch value. \\
\addlinespace
\pkg{pdflscape} &
Landscape pages with PDF rotation metadata. &
Use for genuinely wide tables or figures. Inspect orientation, headers, and
printed output. \\
\addlinespace
\pkg{multicol} &
Balanced multiple text columns. &
Suitable for compact reference material; awkward for wide mathematics,
figures, reading order, and accessibility. \\
\addlinespace
\pkg{enumitem} &
List labels and spacing. &
Define restrained reusable list styles. Excessively compressed lists reduce
readability and may violate a class design. \\
\addlinespace
\pkg{listings} &
Typeset source code without shell execution. &
Configure language, wrapping, escape behaviour, and captions. It is a
typesetting tool, not a syntax validator. \\
\addlinespace
\pkg{minted} &
Syntax highlighting through Pygments. &
Can provide richer highlighting but normally requires external execution and a
Python dependency. Use only when the build and security policy permit it. \\
\addlinespace
\pkg{tcolorbox} &
Breakable pedagogical or technical boxes. &
Offers many optional libraries and dependencies. Keep visual hierarchy modest
and verify page breaks in long boxes. \\
\addlinespace
\pkg{csquotes} &
Context-sensitive quotation commands and bibliography integration. &
Recommended with many \pkg{biblatex} language workflows. It does not decide
whether a quotation is licensed or accurately attributed. \\
\addlinespace
\pkg{biblatex} &
Citation and bibliography management in \LaTeX{}. &
Select a style and backend deliberately. Publisher workflows may require
BibTeX, a generated \file{.bbl}, or their own system instead. \\
\addlinespace
Biber &
Bibliography backend commonly used with \pkg{biblatex}. &
It is a separate program. Run it through \pkg{latexmk} or the documented
sequence; do not substitute BibTeX when the project specifies Biber. \\
\addlinespace
\pkg{natbib} &
Author--year and numeric citation commands for BibTeX-era styles. &
Widely used by publisher classes. Do not load alongside \pkg{biblatex} unless
explicit compatibility documentation supports the setup. \\
\addlinespace
\pkg{hyperref} &
Links, anchors, bookmarks, and PDF metadata. &
Usually load late. Supply plain bookmark text for complex headings and audit
anonymous metadata. \\
\addlinespace
\pkg{cleveref} &
Type-aware cross-references such as Figure 3. &
Load after \pkg{hyperref}; configure names for custom environments and preserve
semantic labels. \\
\addlinespace
\pkg{varioref} &
Contextual page-aware references. &
Useful in some long works; automatic phrases can become awkward after
pagination changes, so proofread the final text. \\
\addlinespace
\pkg{makeidx} &
Basic index collection and printing. &
Requires an indexing pass. Index terms editorially rather than marking every
occurrence. \\
\addlinespace
\pkg{imakeidx} &
Enhanced index management. &
Supports multiple indexes and build integration; may require shell or build
configuration depending on the workflow. \\
\addlinespace
\pkg{glossaries} &
Generated glossaries, acronym lists, and symbols. &
Powerful for large projects but adds database and build complexity. Document
the glossary tool and define terms at first use in prose. \\
\addlinespace
\pkg{acronym} &
Focused acronym definition and use. &
Can suit a moderate acronym list. Compare with \pkg{glossaries} and a manual
list rather than loading both without a plan. \\
\addlinespace
\pkg{mhchem} &
Chemical formulae and reaction notation. &
Use \cmd{ce} according to the current manual. Verify conventions for isotopes,
charges, states, equilibrium, and specialist chemistry. \\
\addlinespace
\pkg{chemfig} &
Programmatic chemical structure diagrams. &
Useful for editable structures; complex source and disciplinary validation are
still required. \\
\addlinespace
\pkg{algorithmicx} with \pkg{algpseudocode} &
Structured pseudocode. &
Define inputs, outputs, and semantics in prose. Do not confuse an algorithm
float, its pseudocode language, and executable validation. \\
\addlinespace
\pkg{tikz} and PGF &
Programmatic vector diagrams and plots. &
Excellent for reproducible simple diagrams. Large graphics can slow builds;
externalise only with a documented toolchain and preserve original source. \\
\addlinespace
\pkg{pgfplots} &
Data-driven scientific plots using PGF/TikZ. &
Set a compatibility version, keep source data, and audit axes and uncertainty.
For large data, generate a vector figure in analysis software instead. \\
\addlinespace
\pkg{standalone} &
Compile individual figures or snippets independently. &
Useful for original TikZ figures. Record the source-to-PDF build and include
only the derived PDF in workflows that prohibit shell execution. \\
\addlinespace
\pkg{lineno} &
Review line numbers. &
Inspect behaviour around equations, floats, and two-column layouts. A publisher
class may supply its own line-number option. \\
\addlinespace
\pkg{changes} and \pkg{latexdiff} &
Visible revision markup and source comparison. &
Review generated markup carefully, especially commands and moved text. Remove
change apparatus from a clean accepted release unless requested. \\
\addlinespace
\pkg{latexmk} &
Dependency-aware build automation. &
Configure the main engine, bibliography, index, glossary, and output directory
in a project-level file. A successful run still needs log and PDF inspection. \\
\bottomrule
\end{longtable}

\section*{Package-selection checklist}

\begin{enumerate}
  \item Name the missing capability.
  \item Check whether the class or kernel already provides it.
  \item Read the current package introduction, minimal example, options, and
        compatibility notes.
  \item Add the package once with the smallest needed configuration.
  \item Compile an MWE and the complete document.
  \item Record extra programs, permissions, or security requirements.
  \item Remove the package if its capability is no longer used.
\end{enumerate}
